% ============================================================================
% CAPÍTULO 2: FERRAMENTAS COMPUTACIONAIS PARA BIOLOGIA DE SISTEMAS
% Bioinformática para Biologia de Sistemas
% ============================================================================

% ============================================================================
% TEMPLATE BASE PARA APRESENTAÇÕES EM BEAMER
% Bioinformática para Biologia de Sistemas
% ============================================================================

\documentclass[12pt, aspectratio=169]{beamer}

% ============================================================================
% PACOTES ESSENCIAIS
% ============================================================================
\usepackage[utf8]{inputenc}
\usepackage[T1]{fontenc}
\usepackage[brazilian]{babel}
\usepackage{amsmath, amssymb, amsthm}
\usepackage{graphicx}
\usepackage{xcolor}
\usepackage{tikz}
\usetikzlibrary{arrows.meta, shapes, positioning, calc, decorations.pathreplacing, decorations.shapes, decorations.pathmorphing}
\usepackage{listings}
\usepackage{booktabs}
\usepackage{multirow}
\usepackage{hyperref}
\usepackage{chemfig}
\usepackage{chemarr}

% ============================================================================
% CONFIGURAÇÃO DO TEMA
% ============================================================================
\usetheme{Madrid}
\usecolortheme{default}

% Cores personalizadas
\definecolor{azulescuro}{RGB}{0, 51, 102}
\definecolor{azulclaro}{RGB}{51, 153, 255}
\definecolor{verdeacido}{RGB}{102, 204, 0}
\definecolor{laranjacel}{RGB}{255, 153, 51}
\definecolor{roxodna}{RGB}{153, 51, 255}

\setbeamercolor{structure}{fg=azulescuro}
\setbeamercolor{palette primary}{bg=azulescuro,fg=white}
\setbeamercolor{palette secondary}{bg=azulclaro,fg=white}
\setbeamercolor{palette tertiary}{bg=azulescuro,fg=white}
\setbeamercolor{palette quaternary}{bg=azulclaro,fg=white}
\setbeamercolor{block title}{bg=azulescuro,fg=white}
\setbeamercolor{block body}{bg=azulclaro!10,fg=black}
\setbeamercolor{block title alerted}{bg=red!80,fg=white}
\setbeamercolor{block body alerted}{bg=red!10,fg=black}
\setbeamercolor{block title example}{bg=verdeacido!80,fg=white}
\setbeamercolor{block body example}{bg=verdeacido!10,fg=black}

% ============================================================================
% CONFIGURAÇÃO DE CÓDIGO
% ============================================================================
\lstset{
    language=Python,
    basicstyle=\ttfamily\small,
    keywordstyle=\color{azulescuro}\bfseries,
    commentstyle=\color{verdeacido},
    stringstyle=\color{laranjacel},
    numbers=left,
    numberstyle=\tiny\color{gray},
    stepnumber=1,
    numbersep=5pt,
    backgroundcolor=\color{white},
    showspaces=false,
    showstringspaces=false,
    showtabs=false,
    frame=single,
    rulecolor=\color{black},
    tabsize=2,
    captionpos=b,
    breaklines=true,
    breakatwhitespace=false,
    escapeinside={\%*}{*)},
    xleftmargin=10pt,
    xrightmargin=10pt
}

% Definir estilo para R
\lstdefinestyle{Rstyle}{
    language=R,
    basicstyle=\ttfamily\small,
    keywordstyle=\color{azulescuro}\bfseries,
    commentstyle=\color{verdeacido},
    stringstyle=\color{laranjacel}
}

% Definir estilo para MATLAB
\lstdefinestyle{MATLABstyle}{
    language=Matlab,
    basicstyle=\ttfamily\small,
    keywordstyle=\color{azulescuro}\bfseries,
    commentstyle=\color{verdeacido},
    stringstyle=\color{laranjacel}
}

% ============================================================================
% COMANDOS PERSONALIZADOS
% ============================================================================

% Comando para equações destacadas
\newcommand{\destaque}[1]{%
    \begin{alertblock}{}
        \begin{center}
            #1
        \end{center}
    \end{alertblock}
}

% Comando para definições
\newcommand{\definicao}[2]{%
    \begin{block}{Definição: #1}
        #2
    \end{block}
}

% Comando para exemplos
\newcommand{\exemplo}[2]{%
    \begin{exampleblock}{Exemplo: #1}
        #2
    \end{exampleblock}
}

% Comando para observações importantes
\newcommand{\importante}[1]{%
    \begin{alertblock}{Importante!}
        #1
    \end{alertblock}
}

% Comandos matemáticos úteis
\newcommand{\R}{\mathbb{R}}
\newcommand{\N}{\mathbb{N}}
\newcommand{\Z}{\mathbb{Z}}
\newcommand{\dif}{\mathrm{d}}
\newcommand{\parcial}[2]{\frac{\partial #1}{\partial #2}}
\newcommand{\derivada}[2]{\frac{\dif #1}{\dif #2}}
\newcommand{\vect}[1]{\boldsymbol{#1}}

% Comandos biológicos
\newcommand{\gene}[1]{\textit{#1}}
\newcommand{\proteina}[1]{\textsc{#1}}
\newcommand{\especie}[1]{\textit{#1}}

% ============================================================================
% CONFIGURAÇÃO DE RODAPÉ
% ============================================================================
\setbeamertemplate{footline}{
  \leavevmode%
  \hbox{%
  \begin{beamercolorbox}[wd=.33\paperwidth,ht=2.25ex,dp=1ex,center]{author in head/foot}%
    \usebeamerfont{author in head/foot}\insertshortauthor
  \end{beamercolorbox}%
  \begin{beamercolorbox}[wd=.34\paperwidth,ht=2.25ex,dp=1ex,center]{title in head/foot}%
    \usebeamerfont{title in head/foot}\insertshorttitle
  \end{beamercolorbox}%
  \begin{beamercolorbox}[wd=.33\paperwidth,ht=2.25ex,dp=1ex,right]{date in head/foot}%
    \usebeamerfont{date in head/foot}\insertshortdate{}\hspace*{2em}
    \insertframenumber{} / \inserttotalframenumber\hspace*{2ex}
  \end{beamercolorbox}}%
  \vskip0pt%
}

% ============================================================================
% CONFIGURAÇÃO DE NAVEGAÇÃO
% ============================================================================
\setbeamertemplate{navigation symbols}{}

% ============================================================================
% CONFIGURAÇÃO DE ITEMIZE/ENUMERATE
% ============================================================================
\setbeamertemplate{itemize items}[circle]
\setbeamertemplate{enumerate items}[default]

% ============================================================================
% FIM DO TEMPLATE
% ============================================================================


\title[Cap. 2: Ferramentas Computacionais]{Capítulo 2: Ferramentas Computacionais para Biologia de Sistemas}
\subtitle{Bioinformática para Biologia de Sistemas}
\author{Ney Lemke}
\institute{DBF-Unesp}
\date{\today}

\begin{document}

% ============================================================================
% SLIDE DE TÍTULO
% ============================================================================
\begin{frame}
\titlepage
\end{frame}

% ============================================================================
% SUMÁRIO
% ============================================================================
\begin{frame}{Sumário}
\tableofcontents
\end{frame}

% ============================================================================
% SEÇÃO 1: HARDWARE E ARQUITETURAS
% ============================================================================
\section{Hardware: o que há dentro do computador}

\begin{frame}{Por que um biólogo precisa entender computadores?}
\begin{block}{Objetivos de Aprendizagem}
\begin{itemize}
    \item Entender os componentes fundamentais de um computador
    \item Conhecer as arquiteturas disponíveis hoje para ciência
    \item Saber escolher a ferramenta certa para cada problema
    \item Dominar o fluxo de trabalho científico reprodutível
\end{itemize}
\end{block}

\vspace{0.3cm}

\importante{
Você não precisa ser programador para fazer ciência computacional — mas precisa entender o \textbf{vocabulário} e as \textbf{ferramentas} do ambiente digital.
}
\end{frame}

\begin{frame}[shrink=15]{Os componentes fundamentais}
\begin{columns}
\column{0.5\textwidth}
\begin{block}{CPU — Unidade Central de Processamento}
\begin{itemize}
    \item O ``cérebro'' do computador
    \item Executa instruções sequencialmente
    \item 4–64 núcleos em desktops modernos
    \item Ótima para tarefas complexas e variadas
\end{itemize}
\end{block}
\vspace{0.2cm}
\begin{block}{RAM — Memória de Acesso Aleatório}
\begin{itemize}
    \item Área de trabalho temporária
    \item Dados desaparecem ao desligar
    \item 8–256 GB em estações de trabalho
\end{itemize}
\end{block}

\column{0.5\textwidth}
\begin{block}{Armazenamento (SSD/HDD)}
\begin{itemize}
    \item Dados persistentes
    \item SSD: rápido (GB/s), caro
    \item HDD: lento, barato, grande capacidade
\end{itemize}
\end{block}
\vspace{0.2cm}
\begin{exampleblock}{Analogia Biológica}
CPU = ribossomo\\
RAM = citoplasma (área de trabalho)\\
SSD = DNA (memória permanente)
\end{exampleblock}
\end{columns}
\end{frame}

\begin{frame}[shrink=5]{GPU: processamento paralelo massivo}
\definicao{GPU — Graphics Processing Unit}{
Processador com milhares de núcleos simples, originalmente criado para gráficos, hoje essencial para Machine Learning, simulações moleculares e análise de imagens biomédicas.
}

\vspace{0.3cm}

\begin{columns}
\column{0.5\textwidth}
\textbf{CPU vs GPU:}
\begin{center}
\begin{tabular}{lcc}
\toprule
& \textbf{CPU} & \textbf{GPU} \\
\midrule
Núcleos & 4--64 & 3\,000--16\,000 \\
Velocidade/núcleo & Alta & Baixa \\
Paralelismo & Baixo & Altíssimo \\
Memória & GBs RAM & GBs VRAM \\
\bottomrule
\end{tabular}
\end{center}

\column{0.48\textwidth}
\begin{exampleblock}{Uso em Biologia}
\begin{itemize}
    \item AlphaFold2 (dobramento de proteínas)
    \item Redes neurais para imagens de microscopia
    \item Dinâmica molecular (GROMACS, NAMD)
    \item Sequenciamento por deep learning
\end{itemize}
\end{exampleblock}
\end{columns}
\end{frame}

\begin{frame}{Arquiteturas de computação disponíveis hoje}
\begin{columns}
\column{0.48\textwidth}
\begin{block}{Laptop / Desktop}
\begin{itemize}
    \item Uso pessoal, análises pequenas
    \item Python, R, Jupyter local
    \item Limitado por RAM e CPU
\end{itemize}
\end{block}
\vspace{0.2cm}
\begin{block}{Workstation}
\begin{itemize}
    \item 32--512 GB RAM, GPU dedicada
    \item Análises de tamanho médio
    \item Compartilhado no laboratório
\end{itemize}
\end{block}

\column{0.48\textwidth}
\begin{block}{Cluster / HPC}
\begin{itemize}
    \item Centenas a milhares de CPUs
    \item Jobs enviados via fila (SLURM)
    \item Usado para simulações longas
\end{itemize}
\end{block}
\vspace{0.2cm}
\begin{block}{Cloud (AWS, GCP, Azure)}
\begin{itemize}
    \item Escala sob demanda
    \item Paga pelo uso (por hora)
    \item GPUs disponíveis instantaneamente
\end{itemize}
\end{block}
\end{columns}
\end{frame}

\begin{frame}{Qual arquitetura usar?}
\begin{center}
\begin{tabular}{lll}
\toprule
\textbf{Tarefa} & \textbf{Arquitetura recomendada} & \textbf{Exemplo} \\
\midrule
Aprender Python/R & Laptop / Google Colab & Análise exploratória \\
Análise de RNA-seq & Workstation ou cluster & DESeq2, edgeR \\
Dobramento de proteína & GPU (Colab Pro / HPC) & AlphaFold2 \\
Genômica populacional & HPC (cluster) & GATK, VCF pipeline \\
Simulação de rede & Workstation & NetworkX, igraph \\
Deploy de modelo ML & Cloud (AWS/GCP) & API de predição \\
\bottomrule
\end{tabular}
\end{center}

\vspace{0.2cm}
\importante{
Comece sempre no ambiente mais simples. Só migre para HPC/cloud quando o problema \textbf{realmente} exigir.
}
\end{frame}

% ============================================================================
% SEÇÃO 2: SISTEMAS OPERACIONAIS E SOFTWARE LIVRE
% ============================================================================
\section{Sistemas Operacionais e Software Livre}

\begin{frame}[shrink=10]{O que é um Sistema Operacional?}
\definicao{Sistema Operacional (SO)}{
Camada de software que gerencia o hardware do computador e fornece serviços para os programas. É o intermediário entre você, seus programas e o processador, a memória e o disco.
}

\vspace{0.3cm}

\begin{columns}
\column{0.5\textwidth}
\textbf{O que o SO faz:}
\begin{itemize}
    \item Gerencia processos (quem usa a CPU)
    \item Controla memória RAM
    \item Organiza o sistema de arquivos
    \item Gerencia dispositivos (teclado, rede, GPU)
    \item Fornece interface ao usuário (gráfica ou terminal)
\end{itemize}

\column{0.48\textwidth}
\begin{exampleblock}{Analogia}
O SO é como o sistema nervoso autônomo do computador: você não precisa pensar em como o coração bate (CPU executa instruções) — ele cuida disso enquanto você foca no que importa.
\end{exampleblock}
\end{columns}
\end{frame}

\begin{frame}[shrink=10]{Windows, Linux e macOS: as três famílias}
\begin{center}
\begin{tabular}{lccc}
\toprule
& \textbf{Windows} & \textbf{Linux} & \textbf{macOS} \\
\midrule
Criador         & Microsoft   & Comunidade (Linus Torvalds) & Apple \\
Custo           & Pago        & Gratuito        & Incluso no Mac \\
Código-fonte    & Fechado     & Aberto          & Parcialmente aberto \\
Uso em servidores & Raro     & Dominante (90\%+) & Raro \\
Uso em desktops & Dominante  & Crescente       & $\sim$15\% \\
Terminal padrão & PowerShell  & bash/zsh        & zsh \\
Instalação de pacotes & .exe  & apt/yum/conda  & brew/conda \\
\bottomrule
\end{tabular}
\end{center}

\vspace{0.2cm}
\importante{
Servidores, clusters e nuvem rodam \textbf{Linux}. Saber o básico de Linux é indispensável para bioinformática.
}
\end{frame}

\begin{frame}[shrink=10]{Linux: distribuições mais usadas em ciência}
\begin{columns}
\column{0.5\textwidth}
\begin{block}{O que é uma distribuição Linux?}
O Linux é apenas o \textbf{kernel} (núcleo). Uma distribuição empacota o kernel com ferramentas, gerenciador de pacotes e interface gráfica. Existem centenas — cada uma com foco diferente.
\end{block}

\vspace{0.2cm}
\begin{exampleblock}{Mais usadas em bioinformática}
\begin{itemize}
    \item \textbf{Ubuntu} — mais fácil para iniciantes
    \item \textbf{Debian} — estável, base do Ubuntu
    \item \textbf{CentOS / Rocky Linux} — padrão em clusters HPC
    \item \textbf{Fedora} — atualizado, bom para workstations
\end{itemize}
\end{exampleblock}

\column{0.48\textwidth}
\begin{block}{Windows com Linux: WSL}
O \textbf{Windows Subsystem for Linux (WSL)} permite rodar Ubuntu dentro do Windows sem dual-boot — boa opção para quem precisa de Windows mas quer usar ferramentas Linux.
\end{block}

\vspace{0.2cm}
\begin{alertblock}{Dica prática}
Se você usa Windows, instale o WSL2 + Ubuntu. Se usa Mac, o terminal já é Unix-compatível — você está a um passo do ambiente científico.
\end{alertblock}
\end{columns}
\end{frame}

\begin{frame}[shrink=10]{Software Livre vs. Software Proprietário}
\definicao{Software Livre (Free/Open Source Software — FOSS)}{
Software cujo código-fonte é público, podendo ser estudado, modificado e redistribuído por qualquer pessoa, respeitando a licença. ``Free'' significa \textbf{liberdade}, não necessariamente gratuito.
}

\vspace{0.2cm}

\begin{columns}
\column{0.5\textwidth}
\begin{exampleblock}{Software Livre em biologia}
\begin{itemize}
    \item Python, R, Julia — linguagens
    \item Linux, Ubuntu — sistemas operacionais
    \item BLAST, BWA, GATK — ferramentas bioinformáticas
    \item Bioconductor, BioPython — bibliotecas
    \item Galaxy, Snakemake — plataformas de análise
\end{itemize}
\end{exampleblock}

\column{0.48\textwidth}
\begin{alertblock}{Software Proprietário em biologia}
\begin{itemize}
    \item Microsoft Office, Windows
    \item MATLAB (licença cara)
    \item CLC Genomics Workbench
    \item Geneious (análise de sequências)
    \item GraphPad Prism (estatística)
\end{itemize}
\end{alertblock}
\end{columns}
\end{frame}

\begin{frame}[shrink=15]{Por que Software Livre importa para a ciência?}
\begin{columns}
\column{0.5\textwidth}
\begin{block}{Vantagens científicas do FOSS}
\begin{itemize}
    \item \textbf{Reprodutibilidade}: qualquer pesquisador pode inspecionar o código
    \item \textbf{Auditabilidade}: algoritmos podem ser verificados e corrigidos
    \item \textbf{Sem vendor lock-in}: você não perde seus dados se a empresa fechar
    \item \textbf{Custo zero}: acessível a pesquisadores de qualquer país
    \item \textbf{Colaboração global}: comunidades como Bioconductor têm milhares de contribuidores
\end{itemize}
\end{block}

\column{0.48\textwidth}
\begin{block}{Tipos de licença comuns}
\begin{itemize}
    \item \textbf{MIT / BSD} — use como quiser, inclusive em software proprietário
    \item \textbf{GPL} — derivados também devem ser livres (``copyleft'')
    \item \textbf{Apache 2.0} — permissiva, com proteção de patentes
    \item \textbf{CC BY} — para dados e artigos (Creative Commons)
\end{itemize}
\end{block}

\exemplo{Impacto real}{
O Linux, Python, R e toda a pilha de bioinformática moderna são software livre. A ciência aberta depende de ferramentas abertas.
}
\end{columns}
\end{frame}

% ============================================================================
% SEÇÃO 3: TERMINAL E LINHA DE COMANDO
% ============================================================================
\section{Terminal e linha de comando}

\begin{frame}[fragile]{O terminal: por que não fugir dele}
\definicao{Terminal / Shell}{
Interface de texto para interagir com o sistema operacional. No Linux/Mac é o \texttt{bash} ou \texttt{zsh}; no Windows é o PowerShell ou WSL (Windows Subsystem for Linux).
}

\vspace{0.3cm}

\begin{columns}
\column{0.5\textwidth}
\textbf{Por que usar o terminal?}
\begin{itemize}
    \item Automatiza tarefas repetitivas
    \item Necessário para usar clusters e Docker
    \item Permite trabalhar em servidores remotos
    \item Reprodutibilidade: comandos são registráveis
\end{itemize}

\column{0.48\textwidth}
\begin{exampleblock}{Analogia}
O terminal é como o pipetador automático da computação: a primeira vez parece complicado, mas depois de aprender, você nunca mais quer fazer na mão.
\end{exampleblock}
\end{columns}
\end{frame}

\begin{frame}[fragile, shrink=20]{Comandos essenciais do terminal}
\begin{columns}
\column{0.5\textwidth}
\begin{lstlisting}[language=bash, caption=Navegação e arquivos, basicstyle=\ttfamily\footnotesize]
# Ver onde estou
pwd

# Listar arquivos
ls -lh

# Entrar em uma pasta
cd minha-analise/

# Criar pasta
mkdir resultados

# Copiar arquivo
cp dados.csv resultados/

# Mover/renomear
mv draft.py analise.py

# Ver conteudo de arquivo
cat sequencias.fasta | head -20
\end{lstlisting}

\column{0.48\textwidth}
\begin{lstlisting}[language=bash, caption=Busca e processos, basicstyle=\ttfamily\footnotesize]
# Buscar texto em arquivo
grep "ATGC" sequencias.fasta

# Contar linhas
wc -l dados.csv

# Ver processos rodando
top

# Rodar script Python
python3 analise.py

# Ver historico de comandos
history | grep python
\end{lstlisting}
\end{columns}
\end{frame}

\begin{frame}[fragile, shrink=15]{Conexão remota com SSH}
\definicao{SSH — Secure Shell}{
Protocolo para acessar computadores remotos (clusters, servidores) com segurança, como se você estivesse na frente do terminal deles.
}

\begin{lstlisting}[language=bash, basicstyle=\ttfamily\footnotesize]
# Conectar a um servidor remoto
ssh usuario@cluster.unesp.br

# Copiar arquivos para o servidor
scp meu_script.py usuario@cluster.unesp.br:~/analises/

# Copiar resultados de volta
scp usuario@cluster.unesp.br:~/resultados/output.csv .

# Transferencia de pastas inteiras
scp -r pasta_local/ usuario@servidor:~/destino/
\end{lstlisting}

\vspace{0.2cm}
\begin{exampleblock}{Dica prática}
Configure chaves SSH (\texttt{ssh-keygen}) para não precisar digitar senha toda vez.
\end{exampleblock}
\end{frame}

% ============================================================================
% SEÇÃO 3: LINGUAGENS DE PROGRAMAÇÃO
% ============================================================================
\section{Linguagens de programação}

\begin{frame}[shrink=15]{O que é uma linguagem de programação?}
\definicao{Linguagem de Programação}{
Um conjunto de regras e vocabulário que permite dar instruções precisas a um computador. É uma tradução entre a lógica humana e as operações binárias da máquina.
}

\vspace{0.3cm}

\begin{block}{Como funciona?}
\begin{enumerate}
    \item Você escreve código em texto legível (ex: Python)
    \item Um \textbf{interpretador} ou \textbf{compilador} traduz para linguagem de máquina
    \item O processador executa as instruções
\end{enumerate}
\end{block}

\vspace{0.2cm}
\exemplo{Analogia}{
Uma linguagem de programação é como um protocolo de laboratório: você descreve cada passo com precisão e o ``assistente'' (computador) executa exatamente o que foi escrito — sem improviso.
}
\end{frame}

\begin{frame}[fragile, shrink=20]{Python: a linguagem do cientista moderno}
\begin{columns}
\column{0.52\textwidth}
\textbf{Por que Python?}
\begin{itemize}
    \item Sintaxe próxima ao inglês — fácil de aprender
    \item Maior ecossistema científico do mundo
    \item Comunidade enorme com tutoriais em tudo
    \item Gratuito e open-source
    \item Integra com R, C++, Java
\end{itemize}

\vspace{0.2cm}
\textbf{Bibliotecas essenciais para biologia:}
\begin{itemize}
    \item \texttt{numpy}, \texttt{pandas} — dados tabulares
    \item \texttt{matplotlib}, \texttt{seaborn} — visualização
    \item \texttt{scipy} — estatística e matemática
    \item \texttt{biopython} — sequências biológicas
    \item \texttt{scanpy} — single-cell RNA-seq
\end{itemize}

\column{0.46\textwidth}
\begin{lstlisting}[language=Python, caption=Exemplo simples, basicstyle=\ttfamily\footnotesize]
import pandas as pd
import matplotlib.pyplot as plt

# Ler dados de expressao genica
dados = pd.read_csv("expressao.csv")

# Ver as primeiras linhas
print(dados.head())

# Histograma de expressao
dados["gene_A"].hist(bins=30)
plt.xlabel("Expressao (TPM)")
plt.ylabel("Frequencia")
plt.show()
\end{lstlisting}
\end{columns}
\end{frame}

\begin{frame}[fragile, shrink=15]{R: o poder da estatística}
\begin{columns}
\column{0.5\textwidth}
\textbf{Quando usar R?}
\begin{itemize}
    \item Análises estatísticas avançadas
    \item Bioconductor: maior coleção de ferramentas bioinformáticas
    \item Publicações que exigem R (DESeq2, edgeR, limma)
    \item Visualizações publicação-ready com \texttt{ggplot2}
\end{itemize}

\vspace{0.2cm}
\begin{block}{Pacotes Bioconductor essenciais}
\begin{itemize}
    \item \texttt{DESeq2} — RNA-seq diferencial
    \item \texttt{edgeR} — contagens de RNA
    \item \texttt{limma} — microarray e RNA-seq
    \item \texttt{clusterProfiler} — enriquecimento
    \item \texttt{Seurat} — single-cell
\end{itemize}
\end{block}

\column{0.48\textwidth}
\begin{lstlisting}[language=R, style=Rstyle, caption=R básico, basicstyle=\ttfamily\footnotesize]
# Instalar pacote
install.packages("ggplot2")
library(ggplot2)

# Ler dados
dados <- read.csv("expressao.csv")

# Grafico de dispersao
ggplot(dados,
  aes(x=condicao_A,
      y=condicao_B)) +
  geom_point(alpha=0.5) +
  geom_smooth(method="lm") +
  theme_minimal()
\end{lstlisting}
\end{columns}
\end{frame}

\begin{frame}[shrink=5]{Python vs R: quando usar cada um?}
\begin{center}
\begin{tabular}{lcc}
\toprule
\textbf{Tarefa} & \textbf{Python} & \textbf{R} \\
\midrule
Machine learning / deep learning & \checkmark\checkmark & $\circ$ \\
RNA-seq diferencial & $\circ$ & \checkmark\checkmark \\
Visualização para publicação & \checkmark & \checkmark\checkmark \\
Automação e pipelines & \checkmark\checkmark & $\circ$ \\
Análise de redes & \checkmark\checkmark & \checkmark \\
Estatística clássica & \checkmark & \checkmark\checkmark \\
Single-cell (scRNA-seq) & \checkmark\checkmark & \checkmark\checkmark \\
Bioinformática geral & \checkmark\checkmark & \checkmark\checkmark \\
\bottomrule
\end{tabular}
\end{center}

\vspace{0.2cm}
\importante{
\checkmark\checkmark = excelente \quad \checkmark = bom \quad $\circ$ = possível, mas não ideal\\
\textbf{Na prática: aprenda Python primeiro, adicione R conforme a necessidade.}
}
\end{frame}

\begin{frame}[fragile, shrink=15]{Outras linguagens relevantes}
\begin{columns}
\column{0.48\textwidth}
\begin{block}{Julia}
\begin{itemize}
    \item Velocidade de C, sintaxe de Python
    \item Crescendo em modelagem matemática
    \item \texttt{BioJulia} para bioinformática
    \item Ideal para simulações numéricas pesadas
\end{itemize}
\end{block}

\vspace{0.2cm}
\begin{block}{Bash / Shell Script}
\begin{itemize}
    \item Automatizar pipelines no terminal
    \item Processar arquivos FASTQ, VCF em lote
    \item Indispensável para HPC
\end{itemize}
\end{block}

\column{0.48\textwidth}
\begin{block}{SQL}
\begin{itemize}
    \item Consultar bancos de dados biológicos
    \item Usado no backend do UCSC, Ensembl
    \item \texttt{SELECT * FROM genes WHERE organism='Homo sapiens'}
\end{itemize}
\end{block}

\vspace{0.2cm}
\begin{block}{C / C++}
\begin{itemize}
    \item Núcleo de ferramentas como BWA, GATK
    \item Não precisa aprender — mas ajuda a entender
\end{itemize}
\end{block}
\end{columns}
\end{frame}

% ============================================================================
% SEÇÃO 4: JUPYTER E GOOGLE COLAB
% ============================================================================
\section{Jupyter e Google Colab}

\begin{frame}[fragile, shrink=15]{Jupyter Notebook: ciência interativa}
\definicao{Jupyter Notebook}{
Ambiente de computação interativa que combina código, texto, equações e visualizações em um único documento. O padrão moderno para ciência de dados reprodutível.
}

\vspace{0.3cm}

\begin{columns}
\column{0.5\textwidth}
\textbf{Vantagens:}
\begin{itemize}
    \item Executa código célula por célula
    \item Visualizações embutidas nos resultados
    \item Mistura código com documentação
    \item Exporta para PDF, HTML, slides
    \item Base do Google Colab
\end{itemize}

\column{0.48\textwidth}
\begin{exampleblock}{Instalação}
\begin{lstlisting}[language=bash, numbers=none, basicstyle=\ttfamily\small]
pip install jupyterlab
jupyter lab
\end{lstlisting}
\end{exampleblock}

\vspace{0.2cm}
\begin{alertblock}{Limitação importante}
Notebooks fora de ordem de execução podem dar resultados errados. Sempre rode \textbf{Kernel $\to$ Restart \& Run All} antes de compartilhar.
\end{alertblock}
\end{columns}
\end{frame}

\begin{frame}[shrink=15]{Google Colab: o laboratório na nuvem}
\begin{block}{O que é o Google Colab?}
Jupyter Notebook hospedado pelo Google, gratuito, sem instalação, com acesso a GPUs e TPUs. Ideal para começar sem configurar nada.
\end{block}

\vspace{0.3cm}

\begin{columns}
\column{0.5\textwidth}
\textbf{Versão gratuita:}
\begin{itemize}
    \item GPU T4 por até 12h/sessão
    \item 12 GB RAM
    \item 15 GB Google Drive integrado
    \item Colaboração em tempo real
\end{itemize}

\column{0.48\textwidth}
\textbf{Como acessar:}
\begin{enumerate}
    \item Acesse \texttt{colab.research.google.com}
    \item Faça login com conta Google
    \item Crie um novo notebook
    \item Habilite GPU: Runtime $\to$ Change runtime type
\end{enumerate}
\end{columns}

\vspace{0.2cm}
\exemplo{Dica}{
Para rodar AlphaFold2, ESMFold e ferramentas de deep learning em biologia, o Colab Pro (\$\,9,99/mês) com GPU A100 é a opção mais acessível para pesquisadores.
}
\end{frame}

% ============================================================================
% SEÇÃO 5: AMBIENTES VIRTUAIS
% ============================================================================
\section{Ambientes virtuais e gerenciamento de pacotes}

\begin{frame}[fragile, shrink=5]{O problema das dependências}
\begin{alertblock}{Cena comum em laboratório}
``O script funcionava no meu computador, mas no servidor ele dá erro na importação do \texttt{pandas}...''
\end{alertblock}

\vspace{0.3cm}

\begin{block}{Por que isso acontece?}
\begin{itemize}
    \item Diferentes versões de pacotes podem ser incompatíveis
    \item Projeto A precisa de \texttt{numpy 1.24}, projeto B precisa de \texttt{numpy 2.0}
    \item Instalar tudo globalmente causa conflitos
\end{itemize}
\end{block}

\vspace{0.2cm}

\definicao{Ambiente Virtual}{
Um diretório isolado com sua própria instalação do Python e pacotes. Cada projeto tem o seu ambiente, evitando conflitos.
}
\end{frame}

\begin{frame}[fragile, shrink=25]{Conda: o gerenciador de ambientes do cientista}
\begin{columns}
\column{0.5\textwidth}
\begin{lstlisting}[language=bash, caption=Criando e usando ambientes, basicstyle=\ttfamily\footnotesize]
# Criar ambiente com Python 3.12
conda create -n meu-projeto \
    python=3.12

# Ativar ambiente
conda activate meu-projeto

# Instalar pacotes
conda install numpy pandas \
    matplotlib scipy biopython

# Listar ambientes
conda env list

# Exportar para reproducao
conda env export > environment.yml

# Recriar em outra maquina
conda env create -f environment.yml
\end{lstlisting}

\column{0.48\textwidth}
\begin{block}{Por que Conda?}
\begin{itemize}
    \item Gerencia Python \textbf{e} dependências não-Python (C libraries, CUDA)
    \item Canal \texttt{bioconda}: 8\,000+ ferramentas bioinformáticas prontas
    \item Instalação simples: \texttt{miniforge} (recomendado)
\end{itemize}
\end{block}

\vspace{0.2cm}
\begin{exampleblock}{Instalar BWA via bioconda}
\begin{lstlisting}[language=bash, numbers=none, basicstyle=\ttfamily\small]
conda install -c bioconda bwa
\end{lstlisting}
\end{exampleblock}
\end{columns}
\end{frame}

\begin{frame}[fragile, shrink=15]{pip e venv: a alternativa leve}
\begin{columns}
\column{0.5\textwidth}
\begin{lstlisting}[language=bash, caption=Usando venv + pip, basicstyle=\ttfamily\footnotesize]
# Criar ambiente virtual
python3 -m venv .venv

# Ativar (Linux/Mac)
source .venv/bin/activate

# Ativar (Windows)
.venv\Scripts\activate

# Instalar pacotes
pip install numpy pandas scipy

# Salvar dependencias
pip freeze > requirements.txt

# Instalar em outro lugar
pip install -r requirements.txt
\end{lstlisting}

\column{0.48\textwidth}
\begin{block}{Quando usar venv + pip?}
\begin{itemize}
    \item Projetos puramente Python
    \item Ambientes mais leves que conda
    \item Integração com PyPI (200k+ pacotes)
\end{itemize}
\end{block}

\vspace{0.3cm}
\begin{center}
\begin{tabular}{lcc}
\toprule
& \textbf{conda} & \textbf{venv} \\
\midrule
Leve & -- & \checkmark \\
Bioconda & \checkmark & -- \\
Não-Python libs & \checkmark & -- \\
Universalidade & \checkmark & \checkmark \\
\bottomrule
\end{tabular}
\end{center}
\end{columns}
\end{frame}

% ============================================================================
% SEÇÃO 6: GITHUB
% ============================================================================
\section{GitHub: controle de versão e colaboração}

\begin{frame}[fragile]{O que é controle de versão?}
\begin{columns}
\column{0.52\textwidth}
\begin{alertblock}{Sem controle de versão}
\begin{itemize}
    \item \texttt{analise\_final.py}
    \item \texttt{analise\_final\_v2.py}
    \item \texttt{analise\_final\_corrigida.py}
    \item \texttt{analise\_USAR\_ESSE.py}
    \item \texttt{analise\_USAR\_ESSE\_2.py}
\end{itemize}
\end{alertblock}

\column{0.46\textwidth}
\begin{exampleblock}{Com Git}
\begin{itemize}
    \item Um único arquivo: \texttt{analise.py}
    \item Histórico completo de todas as versões
    \item Quem mudou o quê e quando
    \item Voltar a qualquer versão anterior
    \item Colaboração sem conflito
\end{itemize}
\end{exampleblock}
\end{columns}

\vspace{0.3cm}
\definicao{Git}{
Sistema de controle de versão distribuído, criado por Linus Torvalds (2005). Registra cada mudança em arquivos ao longo do tempo, como um ``histórico infinito de Ctrl+Z''.
}
\end{frame}

\begin{frame}[fragile, shrink=15]{Comandos Git essenciais}
\begin{columns}
\column{0.5\textwidth}
\begin{lstlisting}[language=bash, caption=Fluxo básico do Git, basicstyle=\ttfamily\footnotesize]
# Iniciar repositorio
git init

# Ver estado atual
git status

# Adicionar arquivo ao stage
git add analise.py

# Salvar versao (commit)
git commit -m "Adiciona analise de expressao"

# Ver historico
git log --oneline

# Voltar a versao anterior
git checkout abc1234
\end{lstlisting}

\column{0.48\textwidth}
\begin{lstlisting}[language=bash, caption=Colaboracao com GitHub, basicstyle=\ttfamily\footnotesize]
# Clonar repositorio existente
git clone https://github.com/
    usuario/projeto.git

# Enviar para GitHub
git push origin main

# Baixar atualizacoes
git pull

# Criar ramo separado
git checkout -b nova-analise

# Unir ramo ao principal
git merge nova-analise
\end{lstlisting}
\end{columns}
\end{frame}

\begin{frame}[shrink=15]{GitHub: muito além do armazenamento de código}
\begin{columns}
\column{0.5\textwidth}
\begin{block}{O que o GitHub oferece?}
\begin{itemize}
    \item Hospedagem gratuita de repositórios
    \item \textbf{Issues}: rastreamento de bugs e tarefas
    \item \textbf{Pull Requests}: revisão de código em equipe
    \item \textbf{Actions}: automação de testes e pipelines
    \item \textbf{Pages}: publicação de sites e relatórios
    \item \textbf{Releases}: versões publicadas de software
\end{itemize}
\end{block}

\column{0.48\textwidth}
\begin{exampleblock}{Boas práticas para cientistas}
\begin{itemize}
    \item Um repositório por projeto/artigo
    \item \texttt{README.md} com instruções claras
    \item \texttt{requirements.txt} ou \texttt{environment.yml}
    \item Dados no repositório ou link para Zenodo/Figshare
    \item DOI para o repositório via Zenodo
\end{itemize}
\end{exampleblock}
\end{columns}

\vspace{0.2cm}
\importante{
Periódicos como \textit{Nature Methods} e \textit{PLOS Computational Biology} exigem ou recomendam fortemente código disponível publicamente no GitHub.
}
\end{frame}

\begin{frame}{GitHub na prática: fluxo de um projeto científico}
\begin{center}
\begin{tikzpicture}[node distance=1.8cm, every node/.style={font=\small}]
    \node[draw, rounded corners, fill=azulescuro!20, minimum width=2.8cm, minimum height=0.8cm] (init) {Criar repositório};
    \node[draw, rounded corners, fill=azulclaro!20, minimum width=2.8cm, minimum height=0.8cm, right=of init] (clone) {Clonar local};
    \node[draw, rounded corners, fill=verdeacido!20, minimum width=2.8cm, minimum height=0.8cm, right=of clone] (work) {Trabalhar e commitar};
    \node[draw, rounded corners, fill=laranjacel!20, minimum width=2.8cm, minimum height=0.8cm, below=of work] (push) {Push para GitHub};
    \node[draw, rounded corners, fill=roxodna!20, minimum width=2.8cm, minimum height=0.8cm, left=of push] (collab) {Colaborador faz Pull Request};
    \node[draw, rounded corners, fill=azulescuro!20, minimum width=2.8cm, minimum height=0.8cm, left=of collab] (review) {Review e merge};

    \draw[->, thick] (init) -- (clone);
    \draw[->, thick] (clone) -- (work);
    \draw[->, thick] (work) -- (push);
    \draw[->, thick] (push) -- (collab);
    \draw[->, thick] (collab) -- (review);
    \draw[->, thick, dashed] (review.north) to[bend left=30] (work.south);
\end{tikzpicture}
\end{center}

\vspace{0.2cm}
\begin{exampleblock}{Para publicação científica}
Ao submeter o artigo: crie uma \textbf{Release} no GitHub e registre no Zenodo para obter um DOI permanente e citável.
\end{exampleblock}
\end{frame}

% ============================================================================
% SEÇÃO 7: DOCKER
% ============================================================================
\section{Docker: reprodutibilidade e ambientes portáteis}

\begin{frame}[shrink=20]{O problema da reprodutibilidade}
\begin{alertblock}{``Works on my machine'' — o pesadelo da ciência computacional}
Você publica o código, o revisor tenta rodar, e nada funciona. Versão diferente do Python, biblioteca atualizada que quebrou a API, dependência que não existe mais.
\end{alertblock}

\vspace{0.3cm}

\begin{block}{A solução: empacotar tudo junto}
Docker permite criar uma \textbf{imagem} que contém:
\begin{itemize}
    \item O sistema operacional base (Ubuntu, Alpine)
    \item Python/R com versões exatas
    \item Todos os pacotes e dependências
    \item O código e os dados (ou links para eles)
    \item As configurações do ambiente
\end{itemize}
\end{block}

\vspace{0.1cm}
\importante{
Se funciona no seu container, funciona em \textbf{qualquer} computador que tenha Docker instalado.
}
\end{frame}

\begin{frame}[fragile, shrink=25]{Docker: conceitos fundamentais}
\begin{columns}
\column{0.5\textwidth}
\definicao{Imagem Docker}{
``Fotografia'' de um ambiente completo. Imutável, compartilhável, versionável. Análoga ao snapshot de uma máquina virtual, mas muito mais leve.
}

\vspace{0.2cm}
\definicao{Container}{
Uma instância em execução de uma imagem. Isolado do sistema hospedeiro, mas compartilha o kernel do sistema operacional.
}

\column{0.48\textwidth}
\begin{exampleblock}{Analogia biológica}
\textbf{Imagem} = receita de um meio de cultura\\
\textbf{Container} = placa com o meio preparado

Você pode preparar mil placas idênticas a partir da mesma receita, em qualquer laboratório do mundo.
\end{exampleblock}
\end{columns}

\vspace{0.3cm}
\begin{center}
\begin{tabular}{lcc}
\toprule
& \textbf{Máquina Virtual} & \textbf{Docker} \\
\midrule
Tamanho típico & 5--50 GB & 100 MB -- 2 GB \\
Tempo de start & Minutos & Segundos \\
Overhead de CPU & Alto & Mínimo \\
Isolamento & Total & Parcial \\
\bottomrule
\end{tabular}
\end{center}
\end{frame}

\begin{frame}[fragile, shrink=25]{Dockerfile: descrevendo seu ambiente}
\begin{lstlisting}[language=bash, caption=Exemplo de Dockerfile para análise bioinformática, basicstyle=\ttfamily\footnotesize]
# Imagem base com Python 3.12
FROM python:3.12-slim

# Metadados
LABEL description="Analise de RNA-seq - Artigo XYZ 2025"

# Instalar dependencias do sistema
RUN apt-get update && apt-get install -y \
    build-essential curl && rm -rf /var/lib/apt/lists/*

# Copiar e instalar dependencias Python
COPY requirements.txt .
RUN pip install --no-cache-dir -r requirements.txt

# Copiar codigo
COPY scripts/ /app/scripts/

# Diretorio de trabalho
WORKDIR /app

# Comando padrao
CMD ["python", "scripts/analise_principal.py"]
\end{lstlisting}
\end{frame}

\begin{frame}[fragile, shrink=20]{Usando Docker na prática}
\begin{columns}
\column{0.5\textwidth}
\begin{lstlisting}[language=bash, caption=Comandos Docker essenciais, basicstyle=\ttfamily\footnotesize]
# Construir imagem
docker build -t minha-analise .

# Rodar container
docker run minha-analise

# Rodar com pasta local montada
docker run -v $(pwd)/dados:/app/dados \
    minha-analise

# Ver containers ativos
docker ps

# Parar container
docker stop CONTAINER_ID

# Baixar imagem do Docker Hub
docker pull jupyter/scipy-notebook
\end{lstlisting}

\column{0.48\textwidth}
\begin{block}{Docker Hub: imagens prontas}
\begin{itemize}
    \item \texttt{jupyter/scipy-notebook} — Jupyter completo
    \item \texttt{rocker/tidyverse} — R + tidyverse
    \item \texttt{biocontainers/bwa} — BWA aligner
    \item \texttt{broadinstitute/gatk} — GATK oficial
    \item \texttt{nfcore/rnaseq} — pipeline completo
\end{itemize}
\end{block}

\vspace{0.2cm}
\begin{exampleblock}{Para publicação}
Publique sua imagem no Docker Hub e inclua o DOI no artigo para garantia de reprodutibilidade permanente.
\end{exampleblock}
\end{columns}
\end{frame}

% ============================================================================
% SEÇÃO 8: WORKFLOWS
% ============================================================================
\section{Workflows: automatizando pipelines}

\begin{frame}[fragile]{O que é um pipeline bioinformático?}
\begin{block}{Pipeline típico de RNA-seq}
\begin{center}
\begin{tikzpicture}[node distance=0.5cm and 0.6cm, every node/.style={font=\tiny},
    box/.style={draw, rounded corners, minimum width=1.6cm, minimum height=0.55cm, text centered}]
    % Linha 1
    \node[box, fill=azulescuro!20] (fastq)  {FASTQ bruto};
    \node[box, fill=azulclaro!20,  right=of fastq]  (qc)    {FastQC};
    \node[box, fill=verdeacido!20, right=of qc]     (trim)  {Trimmomatic};
    \node[box, fill=laranjacel!20, right=of trim]   (align) {STAR/HISAT2};
    % Linha 2 (abaixo de align, seguindo para a esquerda)
    \node[box, fill=roxodna!20,    below=of align]  (count) {featureCounts};
    \node[box, fill=azulescuro!20, left=of count]   (deseq) {DESeq2};
    \node[box, fill=verdeacido!20, left=of deseq]   (enrich){Enriquecimento};
    \node[box, fill=laranjacel!20, left=of enrich]  (viz)   {Visualização};

    \draw[->, thick] (fastq) -- (qc);
    \draw[->, thick] (qc)    -- (trim);
    \draw[->, thick] (trim)  -- (align);
    \draw[->, thick] (align) -- (count);
    \draw[->, thick] (count) -- (deseq);
    \draw[->, thick] (deseq) -- (enrich);
    \draw[->, thick] (enrich)-- (viz);
\end{tikzpicture}
\end{center}
\end{block}

\vspace{0.2cm}
\importante{
Rodar cada etapa manualmente é lento, propenso a erros e impossível de reproduzir. Gerenciadores de workflow automatizam e documentam cada passo.
}
\end{frame}

\begin{frame}[fragile, shrink=20]{Snakemake: workflows em Python}
\definicao{Snakemake}{
Gerenciador de workflows baseado em Python. Define \textbf{regras} que descrevem como gerar arquivos de saída a partir de entradas, e automaticamente determina a ordem de execução.
}

\begin{lstlisting}[language=Python, caption=Snakefile simplificado, basicstyle=\ttfamily\footnotesize]
# Snakefile
rule alinhamento:
    input:
        reads = "dados/{amostra}.fastq.gz",
        genoma = "genoma/GRCh38.fa"
    output:
        bam = "resultados/{amostra}.bam"
    threads: 8
    shell:
        "STAR --runThreadN {threads} "
        "--genomeDir genoma/ "
        "--readFilesIn {input.reads} "
        "--outSAMtype BAM SortedByCoordinate "
        "--outFileNamePrefix resultados/{wildcards.amostra}"
\end{lstlisting}
\end{frame}

\begin{frame}[fragile]{Nextflow e nf-core: pipelines prontos}
\begin{columns}
\column{0.5\textwidth}
\begin{block}{Nextflow}
\begin{itemize}
    \item Linguagem própria (Groovy-based)
    \item Executa em HPC, cloud, Docker, Singularity
    \item Paralelismo automático e escalável
    \item Muito usado em genômica de larga escala
\end{itemize}
\end{block}

\column{0.48\textwidth}
\begin{exampleblock}{nf-core: pipelines prontos}
Coleção de \textbf{pipelines validados} e mantidos pela comunidade:
\begin{itemize}
    \item \texttt{nf-core/rnaseq} — RNA-seq completo
    \item \texttt{nf-core/chipseq} — ChIP-seq
    \item \texttt{nf-core/sarek} — variantes somáticas
    \item \texttt{nf-core/scrnaseq} — single-cell
\end{itemize}
\begin{lstlisting}[language=bash, numbers=none, basicstyle=\ttfamily\tiny]
nextflow run nf-core/rnaseq \
  -profile docker \
  --input samplesheet.csv \
  --genome GRCh38
\end{lstlisting}
\end{exampleblock}
\end{columns}
\end{frame}

% ============================================================================
% SEÇÃO 9: IA COMO COPILOTO
% ============================================================================
\section{IA como copiloto de programação}

\begin{frame}[shrink=5]{LLMs: o novo assistente do cientista}
\begin{block}{O que mudou nos últimos anos}
Modelos de linguagem (ChatGPT, Claude, Gemini) são capazes de escrever, explicar, debugar e adaptar código científico com alta qualidade.
\end{block}

\vspace{0.3cm}

\begin{columns}
\column{0.5\textwidth}
\textbf{O que você pode pedir:}
\begin{itemize}
    \item ``Explique o que este código faz''
    \item ``Converta esse script R para Python''
    \item ``Por que estou recebendo este erro?''
    \item ``Escreva um script para analisar este CSV''
    \item ``Otimize esta função para dados grandes''
\end{itemize}

\column{0.48\textwidth}
\begin{exampleblock}{Ferramentas recomendadas}
\begin{itemize}
    \item \textbf{Claude Code} (Anthropic) — no terminal
    \item \textbf{GitHub Copilot} — no editor de código
    \item \textbf{ChatGPT / Claude.ai} — interface web
    \item \textbf{Gemini no Colab} — integrado ao notebook
\end{itemize}
\end{exampleblock}
\end{columns}
\end{frame}

\begin{frame}{Como usar IA de forma eficaz e responsável}
\begin{columns}
\column{0.5\textwidth}
\begin{block}{Boas práticas}
\begin{itemize}
    \item \textbf{Contexto é tudo}: descreva seus dados, o formato dos arquivos e o objetivo
    \item \textbf{Itere}: peça ajustes em vez de começar do zero
    \item \textbf{Valide sempre}: teste o código antes de usar em análise real
    \item \textbf{Entenda o que gerou}: não publique código que não compreende
\end{itemize}
\end{block}

\column{0.48\textwidth}
\begin{alertblock}{Cuidados importantes}
\begin{itemize}
    \item IA pode gerar código \textbf{plausível mas errado}
    \item Não envie dados sensíveis (pacientes, sequências proprietárias) para IAs externas
    \item Verifique estatísticas — modelos erram em detalhes sutis
    \item Cite ferramentas de IA usadas, quando exigido pelo periódico
\end{itemize}
\end{alertblock}
\end{columns}
\end{frame}

% ============================================================================
% SEÇÃO 10: BOAS PRÁTICAS
% ============================================================================
\section{Boas práticas: reprodutibilidade e ciência aberta}

\begin{frame}{A crise de reprodutibilidade na biologia}
\begin{alertblock}{Dado preocupante}
Um estudo de 2016 na \textit{Nature} mostrou que mais de \textbf{70\% dos pesquisadores} não conseguiram reproduzir experimentos de outros laboratórios, e \textbf{50\%} não conseguiram reproduzir os próprios resultados.
\end{alertblock}

\vspace{0.2cm}

\begin{block}{Causas comuns em estudos computacionais}
\begin{itemize}
    \item Código não disponível ou incompleto
    \item Versões de software não documentadas
    \item Dados pré-processados sem descrever as etapas
    \item Parâmetros ajustados ``a mão'' sem registro
    \item Resultados que dependem de semente aleatória não fixada
\end{itemize}
\end{block}
\end{frame}

\begin{frame}[fragile]{O stack da reprodutibilidade}
\begin{center}
\begin{tikzpicture}[every node/.style={font=\small}]
    \node[draw, fill=azulescuro!30, minimum width=8cm, minimum height=0.7cm, rounded corners] at (0, 3.5) {\textbf{Publicação + DOI} (Zenodo, Figshare)};
    \node[draw, fill=azulclaro!30, minimum width=8cm, minimum height=0.7cm, rounded corners] at (0, 2.6) {\textbf{Repositório Git} (GitHub) — código versionado};
    \node[draw, fill=verdeacido!30, minimum width=8cm, minimum height=0.7cm, rounded corners] at (0, 1.7) {\textbf{Docker} — ambiente fixo e portátil};
    \node[draw, fill=laranjacel!30, minimum width=8cm, minimum height=0.7cm, rounded corners] at (0, 0.8) {\textbf{Conda / venv} — dependências documentadas};
    \node[draw, fill=roxodna!20, minimum width=8cm, minimum height=0.7cm, rounded corners] at (0, 0.0) {\textbf{Jupyter Notebook} — análise documentada passo a passo};
\end{tikzpicture}
\end{center}

\vspace{0.2cm}
\importante{
Reprodutibilidade não é perfeccionismo — é \textbf{ciência funcionando}. Qualquer pessoa deve conseguir chegar ao mesmo resultado com os mesmos dados.
}
\end{frame}

\begin{frame}[fragile]{Checklist para um projeto computacional reprodutível}
\begin{columns}
\column{0.5\textwidth}
\begin{exampleblock}{Estrutura do projeto}
\begin{lstlisting}[numbers=none, basicstyle=\ttfamily\tiny, language=bash]
projeto/
  README.md         # instrucoes
  environment.yml   # conda env
  Dockerfile        # container
  dados/
    brutos/         # originais
    processados/    # derivados
  scripts/
    01_preprocessar.py
    02_analisar.py
    03_visualizar.py
  resultados/
  notebooks/
\end{lstlisting}
\end{exampleblock}

\column{0.48\textwidth}
\begin{block}{Ao finalizar o projeto}
\begin{itemize}
    \item[$\square$] \texttt{README.md} com instruções de execução
    \item[$\square$] Ambiente exportado (\texttt{environment.yml})
    \item[$\square$] Imagem Docker publicada
    \item[$\square$] Repositório GitHub público
    \item[$\square$] Dados em repositório aberto (Zenodo, OSF)
    \item[$\square$] DOI gerado para o código
    \item[$\square$] Sementes aleatórias fixadas no código
\end{itemize}
\end{block}
\end{columns}
\end{frame}

% ============================================================================
% SLIDE FINAL: ROTEIRO DE APRENDIZAGEM
% ============================================================================
\section{Por onde começar?}

\begin{frame}[shrink=15]{Roteiro de aprendizagem para biólogos}
\begin{columns}
\column{0.5\textwidth}
\begin{block}{Mês 1: Fundamentos}
\begin{enumerate}
    \item Google Colab + Python básico
    \item \texttt{pandas} para tabelas de dados
    \item \texttt{matplotlib} para gráficos
    \item Comandos básicos do terminal
\end{enumerate}
\end{block}
\vspace{0.2cm}
\begin{block}{Mês 2--3: Reprodutibilidade}
\begin{enumerate}
    \item Git e GitHub (repositório próprio)
    \item Conda: criar e exportar ambientes
    \item Jupyter Lab local
\end{enumerate}
\end{block}

\column{0.48\textwidth}
\begin{block}{Mês 4--6: Avançar}
\begin{enumerate}
    \item Docker: criar seu primeiro container
    \item Snakemake para automatizar análises
    \item Aplicar ao seu problema de pesquisa
\end{enumerate}
\end{block}
\vspace{0.2cm}
\begin{exampleblock}{Recursos gratuitos}
\begin{itemize}
    \item \textbf{Software Carpentry} — carpentries.org
    \item \textbf{Rosalind} — rosalind.info (bioinformática)
    \item \textbf{Bioconductor workflows} — bioconductor.org
    \item \textbf{The Turing Way} — livro de ciência reprodutível
\end{itemize}
\end{exampleblock}
\end{columns}
\end{frame}

\begin{frame}[shrink=15]{Resumo: o ecossistema computacional do biólogo moderno}
\begin{center}
\begin{tikzpicture}[every node/.style={font=\footnotesize}, node distance=1.2cm]
    % Centro
    \node[draw, circle, fill=azulescuro, text=white, minimum size=1.4cm] (centro) {Seu projeto};

    % Orbita 1
    \node[draw, rounded corners, fill=azulclaro!30, above=of centro] (python) {Python / R};
    \node[draw, rounded corners, fill=verdeacido!30, right=of centro] (github) {GitHub};
    \node[draw, rounded corners, fill=laranjacel!30, below=of centro] (docker) {Docker};
    \node[draw, rounded corners, fill=roxodna!30, left=of centro] (conda) {Conda};

    % Orbita 2
    \node[draw, rounded corners, fill=azulclaro!15, above right=of python] (colab) {Colab/Jupyter};
    \node[draw, rounded corners, fill=verdeacido!15, below right=of github] (snakemake) {Snakemake};
    \node[draw, rounded corners, fill=laranjacel!15, below left=of docker] (zenodo) {Zenodo/DOI};
    \node[draw, rounded corners, fill=roxodna!15, above left=of conda] (hpc) {HPC / Cloud};

    \draw[-] (centro) -- (python);
    \draw[-] (centro) -- (github);
    \draw[-] (centro) -- (docker);
    \draw[-] (centro) -- (conda);
    \draw[-] (python) -- (colab);
    \draw[-] (github) -- (snakemake);
    \draw[-] (docker) -- (zenodo);
    \draw[-] (conda) -- (hpc);
\end{tikzpicture}
\end{center}
\end{frame}

\end{document}
